% Plain math italic, not \mathbf: the symbol is introduced once and never
% reused, so a bold X only reads as an accidental bold word mid-sentence.
Each sample is represented by
$X\in\mathbb{R}^{60\times126}$, where each frame contains
$2\times21\times3=126$ coordinates from two MediaPipe hand slots.
The $x$ and $y$ values are image-normalized, while $z$ is a relative
depth estimate rather than a metric three-dimensional position.
Missing-hand slots are zero-valued. All experiments use the
\texttt{hands126\_v1} processing contract. Table~\ref{tab:profiles}
summarizes the two profiles we evaluate.

\begin{table}[htbp]
\centering
\caption{Evaluated recognition profiles. Performer counts are given as
canonical physical performers over those with complete class coverage.}
\label{tab:profiles}
\small
\setlength{\tabcolsep}{4pt}
\renewcommand{\arraystretch}{1.08}

\begin{tabularx}{\linewidth}{
    >{\raggedright\arraybackslash}p{3.7cm}
    c
    c
    c
    >{\raggedright\arraybackslash}X
}
\toprule
\textbf{Profile}
& \textbf{Classes}
& \textbf{Samples}
& \textbf{Perf.}
& \textbf{Evaluation} \\
\midrule
Hòa Đê workplace signs
& 7 & 445 & 4/4
& Four performer-held-out folds \\
Fingerspelling subset
& 23 & 793 & 5/2
& Two complete-coverage folds \\
\bottomrule
\end{tabularx}
\end{table}

\subsection{Hòa Đê Workplace-Sign Profile}
\label{subsec:hoa_de_profile}

The Hòa Đê profile contains seven workplace signs:
\texttt{rang-muoi}, \texttt{tom}, \texttt{lot-da-ca},
\texttt{lot-vo-tom}, \texttt{cat-ky}, \texttt{danh-vay-ca}, and
\texttt{cat-dau-ca}.

Four adult performers, denoted H01--H04, contributed 445 usable samples,
and all four covered all seven classes. H03 is the originating Deaf
performer. H01, H02, and H04 were instructed performers who had not
habitually used the vocabulary before collection.

Source glosses were entered during capture, with help from a family
member of H03 who communicated with them regularly. We therefore treat
the labels as source annotations grounded in H03's own communication
practice. They are not independent linguistic validation, and we do not
present them as such.

\subsection{Fingerspelling Profile}
\label{subsec:fingerspelling_profile}

The fingerspelling subset contains 793 usable samples across 23 classes:

\begin{equation}
\begin{aligned}
\mathcal{C}_{\mathrm{FS}}=\{&
\mathrm{A},\mathrm{B},\mathrm{C},\mathrm{D},\mathrm{E},
\mathrm{G},\mathrm{H},\mathrm{I},\mathrm{K},\mathrm{L},
\mathrm{M},\mathrm{N},\\
&
\mathrm{O},\mathrm{P},\mathrm{Q},\mathrm{R},\mathrm{S},
\mathrm{T},\mathrm{U},\mathrm{V},\mathrm{X},\mathrm{Y},
\mathrm{Z}\}.
\end{aligned}
\label{eq:fingerspelling_classes}
\end{equation}

This profile is not a complete Vietnamese fingerspelling inventory. It
excludes the Latin letters F, J, and W, as well as the
Vietnamese-specific letters Ă, Â, Đ, Ê, Ô, Ơ, and Ư.

Five physical performers contributed samples, but only F01 and F02
covered all 23 classes. Together, they contributed approximately 89\%
of the usable samples. F01 and F02 correspond to H01 and H02,
respectively; profile-specific pseudonyms are retained to keep the two
collection inventories and split artifacts unambiguous.

One of those performers turned out to have been recorded under two
collection identifiers, \texttt{S005} and \texttt{S007}. We reconciled
the pair to a single physical person in the canonical identity mapping.
Left alone, it would have inflated the profile to six performers where
there are five, and nothing would have stopped a split from placing one
person on both sides of itself.

\subsection{Session and Motion Audits}
\label{subsec:session_motion_audits}

Spreadsheet handling converted several early Hòa Đê session identifiers
into scientific-notation strings, and distinct source identifiers lost
the digits that told them apart. Sample identifiers survived intact, so
no sample is ambiguous, but exact historical session counts cannot be
reconstructed for every record. We therefore describe the collection as
multi-session and make no claim of session-disjoint evaluation.

A descriptive motion audit computed the mean frame-to-frame $\ell_1$
landmark displacement per class. Z, which is dynamic by design, scored
approximately 0.357. Q scored approximately 0.375, higher than Z despite
being intended as a static handshape, and roughly 97\% of its samples
exceeded 0.2. Manual inspection confirmed the cause: the
recordings are contaminated with movement that does not belong to the
sign.

We kept Q in the frozen manifest rather than dropping it, so that the
evaluated task definition stays intact and comparable, and we flag every
fingerspelling result with this caveat. The class is marked for
recollection in a future dataset version.

Raw videos were not retained. Participant-level landmark sequences
remain private, so the conclusions concern the frozen
\texttt{hands126\_v1} representation.